\input{../../../stock/en-tete_v6.tex}

\newcommand{\numerochapitre}{2}
\newcommand{\titrechapitre}{Chapitre \numerochapitre : Logique propositionnelle}

\renewcommand{\headrulewidth}{0.4pt}
\renewcommand{\footrulewidth}{0.4pt}
\fancyfoot[L]{Notes de Raphaël JONTEF}
\fancyfoot[C]{\thepage}
\fancyfoot[R]{Enseirb-Matmeca}
\fancyhead[L]{Cours de Logique et preuve}
\fancyhead[R]{\titrechapitre}

\begin{document}
\input{../../../stock/commands.tex}

% Page de titre custom
\setlength{\headheight}{15pt}
\begin{titlepage}
    \centering
    \vspace*{3cm}
    {\huge Chapitre \numerochapitre\par}
    {\Huge \bfseries\titrechapitre\par}
    \vspace{2cm}
    {\Large Notes rédigées par Raphaël JONTEF\par}
    {\Large avec l'aide de Github Copilot\par}
    \vspace{0.5cm}
    {\normalsize Cours enseigné par Frédéric HERBRETEAU\par}
    \vspace{4cm}
    {\small Enseirb-Matmeca\par}
    {\small filière informatique\par}
    \vfill
    {\large \today\par}
\end{titlepage}

% plan du cours
\tableofcontents
\newpage


% -----------Corps du document--------------

\section{Ambiguïté du langage naturel}

Le language naturel est souvent ambigu~: un même énoncé peut avoir plusieurs interprétations différentes.

\begin{exemple}{}{fromage ou dessert ?}
    Peut-on manger du fromage et du dessert~?
    \begin{itemize}
        \item Oui, si on interprète le "ou" comme un "ou inclusif" (au moins un des deux)
        \item Non, si on interprète le "ou" comme un "ou exclusif" (un seul des deux)
\end{exemple}

\begin{exemple}{}{20/20 chef ?}
    Si on peut résoudre parfaitement n'importe quel exercice du TD pour avoir 20/20, 
    \begin{itemize}
        \item Peut-on avoir 20/20 en faisant n'importe quel exercice du TD ?
        \item Doit-on avoir à faire tous les exercices du TD pour avoir 20/20 ?
    \end{itemize}
\end{exemple}

\section{Introduction à la logique propositionnelle}

En langage naturel, on combine les propositions avec des connecteurs logiques (et, ou, non, si... alors, si et seulement si) pour former des énoncés plus complexes.

\begin{definition}{}{variable propositionnelle}
    Une \notion{variable propositionnelle (ou booléenne)} est une lettre (p, q, r, ...) représentant une proposition à laquelle on associe une valeur de vérité (Vrai ou Faux).
\end{definition}

\begin{definition}{}{notations des connecteurs logiques}
    Les connecteurs de la logique propositionnelle sont notés comme suit~:
    \begin{itemize}
        \item \notion{Négation} : $\neg p$ (non p)
        \item \notion{Conjonction} : $p \land q$ (p et q)
        \item \notion{Disjonction} : $p \lor q$ (p ou q)
        \item \notion{Implication} : $p \implies q$ (si p alors q)
        \item \notion{Équivalence} : $p \iff q$ (p si et seulement si q)
    \end{itemize}
\end{definition}

\begin{proposition}{}{priorités des opérateurs}
    Les priorités des opérateurs logiques sont, de la plus haute à la plus basse~:
    \begin{enumeratebf}
        \item Négation ($\neg$)
        \item Conjonction ($\land$)
        \item Disjonction ($\lor$)
        \item Implication ($\implies$)
        \item Équivalence ($\iff$)
    \end{enumeratebf}
\end{proposition}


\begin{exemple}{}{}
    Les énoncés suivants sont des formules de la logique propositionnelles~:
    \begin{itemize}
        \item $A$
        \item $A \implies B \iff C$
        \item $\neg\neg\neg(A \lor \neg B)$
    \end{itemize}

    Les suivantes n'en sont pas~:
    \begin{itemize}
        \item $A \implies$
        \item $ A + B$
        \item $A \neg B$
\end{exemple}

\begin{definition}{}{table de vérité}
    Une \notion{table de vérité} est un tableau listant toutes les combinaisons possibles des valeurs de vérité des variables propositionnelles d'une formule, ainsi que la valeur de vérité de la formule elle-même pour chaque combinaison.
\end{definition}

\begin{proposition}{}{table de vérité de la négation}
    La table de vérité de la négation est la suivante~:
    \begin{center}
        \begin{tabular}{c|c}
            $p$ & $\neg p$ \\
            \hline
            V & F \\
            F & V \\
        \end{tabular}
    \end{center}
\end{proposition}

\begin{proposition}{}{table de vérité de la conjonction}
    La table de vérité de la conjonction est la suivante~:
    \begin{center}
        \begin{tabular}{c|c|c}
            $p$ & $q$ & $p \land q$ \\
            \hline
            V & V & V \\
            V & F & F \\
            F & V & F \\
            F & F & F \\
        \end{tabular}
    \end{center}
\end{proposition}

\begin{proposition}{}{table de vérité de la disjonction}
    La table de vérité de la disjonction est la suivante~:
    \begin{center}
        \begin{tabular}{c|c|c}
            $p$ & $q$ & $p \lor q$ \\
            \hline
            V & V & V \\
            V & F & V \\
            F & V & V \\
            F & F & F \\
        \end{tabular}
    \end{center}
\end{proposition}

\begin{proposition}{}{table de vérité de l'implication}
    La table de vérité de l'implication est la suivante~:
    \begin{center}
        \begin{tabular}{c|c|c}
            $p$ & $q$ & $p \implies q$ \\
            \hline
            V & V & V \\
            V & F & F \\
            F & V & V \\
            F & F & V \\
        \end{tabular}
    \end{center}
\end{proposition}

\begin{proposition}{}{table de vérité de l'équivalence}
    La table de vérité de l'équivalence est la suivante~:
    \begin{center}
        \begin{tabular}{c|c|c}
            $p$ & $q$ & $p \iff q$ \\
            \hline
            V & V & V \\
            V & F & F \\
            F & V & F \\
            F & F & V \\
        \end{tabular}
    \end{center}
\end{proposition}
    

\begin{definition}{}{valuation}
    Une \notion{valuation} donne une valeur de vérité à chaque variable propositionnelle d'une formule.
\end{definition}

\begin{proposition}{}{généralisation des valuations aux formules}
    Étant donnée une formule $F$ et une valuation $v$, la \notion{valeur de vérité de $F$ sous la valuation $v$} est notée $v(F)$ et est définie de manière inductive~:
    \begin{itemize}
        \item Si $F$ est une variable propositionnelle $p$, alors $v(F) = v(p)$.
        \item Si $F$ est de la forme $\neg G$, alors $v(F) = 1$ si $v(G) = 0$, et $v(F) = 0$ si $v(G) = 1$.
        \item Si $F$ est de la forme $G_1 \land G_2$, alors $v(F) = 1$ si $v(G_1) = 1$ et $v(G_2) = 1$, sinon $v(F) = 0$.
        \item Si $F$ est de la forme $G_1 \lor G_2$, alors $v(F) = 1$ si au moins un des $v(G_i) = 1$, sinon $v(F) = 0$.
        \item Si $F$ est de la forme $G_1 \implies G_2$, alors $v(F) = 0$ si $v(G_1) = 1$ et $v(G_2) = 0$, sinon $v(F) = 1$.
        \item Si $F$ est de la forme $G_1 \iff G_2$, alors $v(F) = 1$ si $v(G_1) = v(G_2)$, sinon $v(F) = 0$.
    \end{itemize}
\end{proposition}

\begin{exemple}{}{}
    Si $v(A) = 0, v(B) = 1, v(C) = 1$, alors $v\Big(\underbrace{(A \implies B)}_1 \land \underbrace{(\neg A \implies C)}_1\Big) = 1$.

\end{exemple}

\begin{remarque}{}{}
    Pour une formule comportant un nombre $n$ de variables propositionnelles, il y a $2^n$ valuations possibles, donc $2^n$ lignes dans la table de vérité.
\end{remarque}

\begin{definition}{}{validité d'une formule}
    Une formule est \notion{valide} si elle est vraie pour toute valuation. On peut également qualifier une telle formule de \notion{tautologie}.
\end{definition}

\begin{remarque}{}{}
    On peut donc montrer qu'une formule est valide en construisant sa table de vérité et en vérifiant que la colonne correspondant à la formule est remplie de Vrai.
\end{remarque}

\begin{definition}{}{notation}
    Dans un système de preuve, une \notion{règle d'inférence} est constituée d'une \notion{famille de prémisses $P_1,\, \dots,\, P_k$}, et d'une \notion{conclusion $C$}. On représente une règle d'inférence par~:
    $$\frac{P_1 \quad \dots \quad P_k}{C}$$
\end{definition}

\begin{remarque}{}{}
    Affirmer que~:
    $$\frac{P_1 \quad \dots \quad P_k}{C}$$
    revient à démontrer que si toutes les prémisses sont vraies, alors la conclusion l'est aussi. Ce qui revient à montrer que~:
    $$(P_1 \land \dots \land P_k) \implies C$$
    est une tautologie.
\end{remarque}

\section{Moyens de preuve en logique propositionnelle}

\subsection{Via table de vérité}
    On peut montrer qu'une formule est valide en construisant sa table de vérité et en vérifiant que la colonne correspondant à la formule est remplie de Vrai.
\subsection{Par preuve sémantique}
    On prend une valuation quelconque, et on montre que la formule est vraie sous cette valuation, peu importe sa définition. C'est un raisonnement en langage naturel.s
\section{Équivalence de formules de la logique propositionnelle}

\begin{definition}{}{équivalence de formules}
    Deux formules $F$ et $G$ sont \notion{logiquement équivalentes} si pour toute valuation $v$, $v(F) = v(G)$. On note alors $F \equiv G$.
\end{definition}

\begin{proposition}{}{lois de De Morgan}
    Les lois de De Morgan sont les suivantes~:
    \begin{itemize}
        \item $\neg (A \land B) \equiv \neg A \lor \neg B$
        \item $\neg (A \lor B) \equiv \neg A \land \neg B$
    \end{itemize}
\end{proposition}

Il existe plusieurs moyens de prouver l'équivalence de deux formules~:
\subsection{Via comparaison des tables de vérité}
    On construit la table de vérité des deux formules, et on vérifie que les colonnes correspondantes sont identiques.
\subsection{Par preuve sémantique}
    On prend une valuation quelconque, et on montre que $v(F) = v(G)$, peu importe la définition de $v$. C'est un raisonnement en langage naturel.
    \begin{exemple}{}{}
        Si $F = A \lor (\neg A \land B)$ est valide, alors $A\lor B$ est vraie. Soit $v$ une validation compatible. Alors $v(F) = 1$.
        \begin{itemize}
            \item Si $v(A) = 1$, alors $v(A \lor B) = 1$.
            \item Si $v(A) = 0$, alors $v(\neg A) = 1$, donc $v(B) = 1$ et donc $v(A \lor B) = 1$.
        \end{itemize}
        On en déduit que pour toute valuation $v$, $v(A \lor B) = 1$, donc $A \lor B$ est valide. Ainsi, $F \equiv A \lor B$.
    \end{exemple}

\section{satisfiabilité de formules de la logique propositionnelle}


\begin{definition}{}{satisfiabilité}
    Une formule $F$ est \notion{satisfiable} s'il existe une valuation $v$ compatible telle que $v(F) = 1$. Dans le cas inverse, on dit que $F$ est \notion{insatisfiable}.
\end{definition}

\begin{theoreme}{}{}
    Une formule $F$ est insatisfiable si et seulement si sa négation $\neg F$ est valide.
\end{theoreme}

\subsection{Problème de satisfiabilité (SAT)}

On définit le problème suivant~:
\code{SAT}~:
\begin{itemize}
    \item \textbf{Entrée} : une formule $F$ de la logique propositionnelle
    \item \textbf{Question} : $F$ est-elle satisfiable, \ie existe-t-il une valuation $v$ telle que $v(F) = 1$~?
\end{itemize}

\begin{remarque}{}{}
    De nombreux problèmes algorithmiques peuvent être réduits au problème SAT, qui est NP-complet.
\end{remarque}

\section{Interlude de décidabilité}

\begin{definition}{}{classe NP}
    La classe NP est l'ensemble des problèmes de décision pour lesquels il existe un \notion{algorithme vérificateur} permettant de \notion{vérifier} qu'une instance est positive à partir d'un couple (instance, certificat) en temps polynomial en la taille de l'instance, avec un certificat de taille polynomiale en fonction de cette taille.
\end{definition}


\begin{definition}{}{réduction polynomiale}
    Soit deux problèmes de décision $\mc{P}_1$ et $\mc{P}_2$ définis par $f_1$ et $f_2$ leurs fonctions de prédicats associées définies sur un $I_1$ et $I_2$ respectivement. On dit que \notion{$\mc{P}_1$ se réduit à $\mc{P}_2$ de façon polynomiale} (noté $\mc{P}_1 \leq_m^\mr{P} \mc{P}_2$) lorsqu'il existe $g:I_1 \to I_2$ calculable avec un algorithme de complexité polynomiale telle que~:
    $$\forall e \in I_1,\, f_1(e) = f_2\Big(g(e)\Big)$$
\end{definition}

\begin{definition}{}{problème de décision NP-difficile}
    Un problème de décision est \notion{NP-difficile} ou \notion{NP-dur} lorsqu'il est possible de se ramener à tout problème de classe NP par une réduction polynomiale depuis ce problème.\\
    Moralement, il est plus difficile que tous les problèmes NP, mais on peut s'y ramener.
\end{definition}

\begin{theoreme}{}{}
    Soit $\mc{P}$ un problème de décision. $\mc{P}$ est NP-difficile si et seulement si~:
    $$\forall \mc{P}' \in \mr{NP},\, \mc{P}' \leq_m^\mr{P} \mc{P}$$
    Si de plus $\mc{P}$ est dans NP, alors $\mc{P}$ est \notion{NP-complet}.
\end{theoreme}


\begin{remarque}{}{}
    Pour montrer qu'un problème $\mc{P}$ est NP-complet~:
    \begin{itemize}
        \item on montre qu'il est de classe NP  l'aide d'un certificat et d'un algorithme vérificateur.
        \item on montre qu'il est NP-difficile en choisissant un problème $\mc{P}_\mr{ref}$ NP-difficile de référence (Comme SAT par exemple), et on montre que $\mc{P}_\mr{ref} \leq_m^\mr{P} \mc{P}$.
    \end{itemize}
\end{remarque}  

\section{Formes normales}

\begin{definition}{}{forme normale disjonctive}
    Une formule est en \notion{forme normale disjonctive (FND)} si elle est une disjonction de clauses conjonctives, où~:
    \begin{itemize}
        \item une \notion{clause conjonctive} est une conjonction de littéraux
        \item un \notion{littéral} est une variable propositionnelle ou sa négation.
    \end{itemize}
\end{definition}

\begin{remarque}{}{}
    La satisfiabilité d'une forme normale disjonctive repose sur celle d'au moins une des clauses conjonctives qui la constituent.
\end{remarque}

\begin{remarque}{}{}
    Bien sûr toute formule de la logique propositionnelle est équivalente à une forme normale disjonctive, mais le problème consistant à trouver cette FND est NP-complet.
\end{remarque}

\begin{definition}{}{forme normale conjonctive}
    Une formule est en \notion{forme normale conjonctive (FNC)} si elle est une conjonction de clauses disjonctives, où~:
    \begin{itemize}
        \item une \notion{clause disjonctive} est une conjonction de littéraux
        \item un \notion{littéral} est une variable propositionnelle ou sa négation.
    \end{itemize}
\end{definition}

\begin{remarque}{}{}
    Là encore, toute formule de la logique propositionnelle est équivalente à une forme normale conjonctive, mais le problème consistant à trouver cette FNC est de classe P. En effet, il existe un algorithme la trouvant en temps polynomial en la taille de la formule d'entrée.
\end{remarque}

    



\end{document}

